% Transcription — fonds Alexandre Grothendieck, shelfmark 138, batch 7 (pages 121-132).
% Established from the facsimile archives/batches/138/batch-07.pdf (cover sheet
% excluded; batch page 1 = archivists' page 121), per the skill
% transcribe-grothendieck.
%
% Pass: Opus 5.5 (claude-opus-5-5), 2026-09-30 — first pass, unchecked against the pages by a human.
%
% Eleven of the twelve pages are transcribed. Page 132 is an unrelated typed
% sheet with nothing in his hand; it is absent.
%
% Contents, in outline:
%   Page 121 — section « 1. »: for a presheaf F on oriented finite maps, the
%     quotient F^♮(C) = F(X)/Aut^0(C) by automorphisms isotopic to the
%     identity, its functoriality, extension to maps with a finite group, and
%     the programme of finding all such F^♮ : finite maps/iso → finite sets.
%   Page 122 — Γ(F^♮) as a colimit over free finite-index subgroups H_α of the
%     cartographic group; example F = P (power set) giving P(S)/Aut(S), the
%     unoriented variant; NB on subsets A of S yielding sub-1-complexes A_X.
%   Page 124 — continues p. 122's last sentence: connectedness conditions a),
%     b), and a choice of vertices, face centres and edge midpoints (1°-3°).
%   Pages 123 and 127 — a separate run: A with its vertices is recovered as
%     φ^{-1}([0,1]) for φ : S → S ramified over 0, 1, ∞ only; p. 127 ends on
%     « Que se passe-t-il dans le cas exceptionnel ? ». Page 123 opens inside
%     a parenthesis whose beginning is not in this batch.
%   Page 125 — Théorème: a finite group acting on a map; equivalent conditions
%     a)-d) (injectivity into isotopy classes, invariant rigidification,
%     realisation as an action on a map), contractibility of the space of
%     invariant rigidifications; start of the proof b) ⇒ a).
%   Page 126 — scratch sheet: Möbius maps z ↦ z/(cz+1-c), c = 0 or 2,
%     g(z) = (2f(z)-1)^2, sketches (some doodles).
%   Pages 128-130 — holomorphic maps: (X, K, S) holomorphic iff there is
%     φ : X → P^1 étale off 0, 1, ∞, ramification 2 over 1, (K, S) =
%     φ^{-1}([0,1], 0); uniqueness of φ; exceptional cases I), II), III)
%     with groups Aff(1,R), C^*, Aff(1,C); conformal representation; fixing
%     φ^{-1}(∞) eliminates them.
%   Page 131 — « Opération » defined by a map (K', S') on S: conditions a)
%     b) c), a') b') c'), d) (a special holomorphic φ : S → S).
%
% Notation fixed for the batch: C = (X, K, S) a map (X the surface, K the
% 1-skeleton, S the sphere/base); his « \mathcal{F} » script F for the functor;
% the sign he sets as an exponent on F (a small hooked mark) is rendered
% \natural and flagged once; S = P^1_C; A, A_x as he writes them.
% His underlinings are set with \emph.
\documentclass[11pt,a4paper]{article}
\input{../preamble/grothendieck}

\folder{138}
\batch{7}
\pages{121}{132}
\foldertitle{Cartes. Etude arithmétique (1976, 77 ?) : bon de commande (s.d.), notes manuscrites (s.d.).}
\dating{[vers 1976-1978]}
\watermark{Édition de démonstration}

\begin{document}

\page{121}
\textbf{1.} \struck{Appelons « \uncertain{opération} » (sur cartes \add{orientées}) toute \ill{} foncteur}
\emph{Soit} \struck{$C = (X, K, S)$ une carte}, soit $\mathcal{F}$ un \struck{faisceau} \add{\uncertain{préfaisceau}} sur
la catégorie des \struck{surfaces} \add{cartes finies \add{orientées}} \struck{compactes, (\ill{} \ill{}
morphismes les morphismes des \ill{}} \struck{(i.e.\ avec ramifications).} [On suppose que
si $(X, H)$ est une surface compacte avec $H$ groupe fini opérant avec ramification isolée, on a
$\mathcal{F}(X/H) = \mathcal{F}(X)^{H}$.]
\marginal{On suppose que $\mathcal{F}$ transforme sommes en produits\ldots}
$\mathrm{Aut}(C)$ opère sur $\mathcal{F}(X)$, donc \uncertain{a fortiori} son sous-groupe
$N(C)$ \add{$= \mathrm{Aut}^{0}(C)$} \struck{des automorphismes} isotopes à l'identité
\uncertain{opère}, soit
\[
  \mathcal{F}^{\natural}(C) = \mathcal{F}(X)/\mathrm{Aut}^{0}(C).
\]
\note{l'exposant de $\mathcal{F}$, ici et dans toute la page, est un petit signe crochu que nous rendons par $\natural$ ; la lecture du signe est incertaine.}
C'est un \struck{foncteur} \add{préfaisceau} en $C$, car si
$C' = (X', K', S') \xrightarrow{\;f\;} (X, K, S)$, \uncertain{on a}
$\mathcal{F}(X) \to \mathcal{F}(X')$ \struck{et \ill{} $\mathrm{Aut}(f)^{0}$}
de façon compatible avec $\mathrm{Aut}(f)^{0}$, donc
\[
  \mathcal{F}(X)/\mathrm{Aut}(f)^{0} \longrightarrow \mathcal{F}(X')/\mathrm{Aut}(f)^{0},
\]
or $\mathrm{Aut}(f)^{0} \xrightarrow{\sim} \mathrm{Aut}(X)^{0}$, d'où par composition
\[
  \mathcal{F}(X)/\mathrm{Aut}(C)^{0} \to \mathcal{F}(X')/\mathrm{Aut}(f)^{0}
  \to \mathcal{F}(X')/\mathrm{Aut}(C')^{0},
\]
i.e.\ $\mathcal{F}^{\natural}(C) \to \mathcal{F}^{\natural}(C')$.
On étend $\mathcal{F}^{\natural}$ en un foncteur sur la catégorie des cartes \uncertain{orientées} $C$ avec
groupe fini d'automorphismes, en posant
\[
  \mathcal{F}^{\natural}(C, G) = \mathcal{F}(C)^{G}/\mathrm{Aut}(C, G)^{0}.
\]
Ceci dit, on se propose de trouver \struck{\ill{}} \uncertain{toutes les} solutions de \struck{$G$}
\[
  \mathcal{F}^{\natural} : \text{Cartes finies}/\text{isomorphismes} \longrightarrow \text{ens.\ finis}
\]
\note{« finies » est ajouté par lui au-dessus de « Cartes ».}

\page{122}
Un argument essentiellement formel montre que
\[
  \Gamma(\mathcal{F}^{\natural}) \simeq \varinjlim_{H_\alpha} \mathcal{F}^{\natural}(C_\alpha, G_\alpha)
\]
\struck{où $C_\alpha$, \ill{}} où $H_\alpha$ parcourt les sous-groupes \emph{\uncertain{libres}} d'indice fini du groupe
\uncertain{cartographique} orienté, $C_\alpha$ est la carte orientée associée à $H_\alpha$, $G_\alpha = \mathcal{C}_2^{+}/H_\alpha$.

\emph{Ex} Soit $\mathcal{F}(X) = \mathfrak{P}(X)$, alors
\[
  \Gamma\mathcal{F}^{\natural} \simeq \mathfrak{P}(S)/\mathrm{Aut}(S),
\]
où $S$ \struck{est} \add{la} sphère \struck{munie d'une} \struck{\ill{} \ill{} $(0,1,\infty)$} \add{munie du segment $[0,1]$ et le point $\infty$,}
qui sert de référence dans notre théorie de \uncertain{rigidification}.

\drawing{une sphère ; sur l'équateur, en pointillé, le point $\infty$ ; en trait gras le segment de $0$ à $1$.}

\emph{NB} Si on avait pris la théorie non orientée, on aurait trouvé
\[
  \Gamma(\mathcal{F}^{\natural}) \simeq \mathfrak{P}(S)^{\underline{\sigma}}/\mathrm{Aut}(S, \underline{\sigma}, 0, 1, \infty)^{0}
\]
avec $\mathrm{Aut}(S, \underline{\sigma}, 0,1,\infty)^{0} \simeq \mathrm{Aut}(S^{+}, 0, 1, \infty)^{0}$
\note{au-dessus de $\mathrm{Aut}$, une accolade porte la mention « $\simeq$ » ; l'exposant souligné $\underline{\sigma}$ est une lecture incertaine (symétrie de conjugaison).}
(où $\underline{\sigma}$ est la conjugaison complexe \uncertain{considérée} comme \uncertain{symétrie} \uncertain{par rapport à} l'« équateur » \add{\uncertain{\ill{}}} qui passe par $0, 1, \infty\ldots$)

\emph{NB} On prouve que, pour que \struck{cette} \add{la} \uncertain{solution} $\Omega$ de $\mathcal{F}^{\natural}$ \add{définie} \struck{associée} à une partie $A$ \add{$\cup$} de
$S$ \struck{définie} \add{associée} \struck{par \ill{} \ill{}} à une carte $C = (X, K, S)$
\struck{des} parties \add{$A_X$} de $X$ qui \struck{sont} un sous-1-complexe \uncertain{topologique}, il faut et il suffit que
la partie de $S$ envisagée soit de ce \uncertain{type}.
Pour que \struck{les} $A_X \cup K$ soit des sous-1-complexes (\uncertain{par} \uncertain{ex.} \uncertain{sous-simpliciaux}) il faut et il suffit
\note{la phrase se poursuit en tête de la page~124 (« que $A \cup [0,1]$ soit un 1-complexe »), non de la page~123.}

\drawing{en marge gauche, un quadrilatère (enveloppe) dont les diagonales et deux segments se rencontrent en un point intérieur.}

\page{123}
\note{la page s'ouvre au milieu d'une parenthèse dont le début ne figure ni à la page~122 ni ailleurs dans ce lot ; elle se poursuit à la page~127.}
(auquel cas les points correspondants des
$A_X$ ne sont pas \add{\uncertain{nécessairement}} des sommets topologiques
\ldots) \struck{on \ill{}} te considérer comme un « milieu
d'arête » des $A_X$.

À l'exclusion de ce cas, on peut dire
que $A$, avec l'ensemble de ses sommets
choisis (qui contient \emph{tous} les sommets \uncertain{qui sont} $\in A$,
topologiques, et les points $0, 1, \infty$ (qu'ils
soient ou non sommets top., et éventuellement
d'autres encore, suivant l'opération
qu'on a \uncertain{définie} i.e.\ les sommets
qu'on attend donner aux $A_X$), et
le choix des arêtes de \uncertain{façon} fait
en conformité avec la condition $1^{\circ}$),
et enfin des deux \emph{\uncertain{orientations}} des
milieux d'arêtes de $A$ (tous distincts
\uncertain{non} de $0, 1, \infty$) est déduit d'une
\uncertain{application} $\varphi : S \to S$ ramifiée au-dessus des
seuls points $0, 1, \infty$, avec ramification
2 au-dessus de $1$, en prenant l'image
inverse de $[0,1]$ --- et l'image inverse
\struck{des} $A_X$ dans $X$ par $f : X \to S$ de
$A$ \uncertain{est} \uncertain{aussi} $(\varphi f)^{-1}([0,1])$, \ldots
$\varphi f$ est \struck{ramifié} un rev.\ \uncertain{ramifié} en dehors de
$0, 1, \infty$ (car $\varphi^{-1}\{0,1,\infty\} \supset \{0,1,\infty\}$ et

\page{124}
que $A \cup [0,1]$ soit un 1-complexe.
Pour que de plus les $A_X$ \struck{soient} \add{définissent} \struck{\ill{}}
sur les $X$ de nouvelles structures de
cartes (i.e.\ $X \setminus A_X$ a comp.\ conn.\ $\simeq$ disques)
il faut et il suffit que a) $A$ \struck{$X$} soit connexe
\struck{a) Si $A$ ne \ill{} \ill{} $0, 1, \infty$, \ill{}}
\struck{b) $A$ contient \ill{} il n'y a \ill{} qui ne soit pas}
\struck{s'il \ill{}} \add{plus des points $0, 1, \infty$ qui sont $\notin A$,}
\add{i.e.\ définit une carte sur $S^{1}$}
\struck{\ill{} \ill{}} \struck{\ill{} \ill{} des autres cartes \ill{} deux}
b) $A$ \struck{\ill{}} ne peut contenir pas plus qu'un seul
des points $0, 1, \infty$.
\note{le passage entre a) et b) est un palimpseste : plusieurs lignes biffées, un cadre, et un ajout encadré dont la place dans la phrase est incertaine.}

\uncertain{Supposons} cette condition \uncertain{satisfaite}. Choisissons
des sommets de $A$, \struck{et} des centres des
faces, des milieux d'arêtes, de façon
à satisfaire à 2 conditions :
1°) Si \struck{\ill{}} des points $0, 1, \infty$ est dans une
face \uncertain{ouverte}, on le choisit comme
centre de cette [il n'y a pas de conflit\ldots]
2°) Si \uncertain{un} [i.e.\ \struck{soit} $s \in A$] on choisit $s$
comme sommet de \struck{$A$} \add{$A$}, \struck{\ill{}} à l'exception
du seul cas qui suit :
\marginal{\uncertain{montrons} que ces \uncertain{choix} \uncertain{obligés} \ill{} \ill{} \ill{} font, \ill{}, les \uncertain{sommets} topolog.\ des $A_X$\ldots}
3°) Si $s = 1 \in A$, et si $s$ est \struck{un} pt
d'ordre $1$ de $A$
i.e.\ un bout

\drawing{un segment issu du point $s$, marqué $A$, se prolongeant par une ligne brisée jusqu'à une boucle.}

\page{125}
\emph{Th.} Soit $C = (X, K, S)$ carte \add{\uncertain{connexe}} \struck{\ill{}}, $G$ un groupe opérant \add{fidèlement} sur $X$
\struck{\uncertain{supposons} que $G \to \mathrm{Aut}_{\mathrm{isot}}(C)$ soit injectif} \uncertain{sur} \uncertain{facilité}.
Conditions équivalentes
\begin{itemize}
  \item[a)] \struck{L'injectivité de $G$} Pour $g \in G$ tel que $g_X$ soit isotope à id
  (en \uncertain{tant} qu'automorphisme de cartes), $g_X = \mathrm{id}_X$, i.e.\ \uncertain{$g = 1$} \uncertain{i.e.}
  $G \to \mathrm{Aut}_{\mathrm{isot}}(C)$ est injectif (i.e.\ $G$ \uncertain{opère} \uncertain{fidèlement} sur $X$ \uncertain{de} $C$ \uncertain{à} \uncertain{isotopie}).
  \item[b)] \struck{\ill{} \ill{} \ill{} \ill{} \ill{} sommet \ill{} de $X$]} \add{\uncertain{Le} \uncertain{stabilisateur} \uncertain{dans} $G$ \uncertain{de} $F_x$ est fini}
  \struck{b) Le noyau de $G \to \mathrm{Aut}_{\mathrm{isot}}(G)$ est fini}
  \item[c)] Il existe une rigidification de $X$ invariante par $G$.
  \item[d)] $(C, G)$ est isomorphe à \struck{la} \uncertain{réalisation} top.\ d'une opération de
  \struck{Dans ces conditions, ou \uncertain{mieux}} $G$ sur une carte \uncertain{\ill{}}.
  \struck{De plus}, l'espace des rigidifications de $X$ invariantes
  par $G$ est contractile.
\end{itemize}
\note{la numérotation des conditions a été remaniée : un b) et un a) sont réécrits à l'encre noire, une flèche en marge ramène la condition sur le noyau de $G \to \mathrm{Aut}_{\mathrm{isot}}$ ; l'ordre ci-dessus suit la dernière main, sous réserve.}

a) $\Rightarrow$ b) $\Rightarrow$ c) évident.

b) $\Rightarrow$ a) Revient à ceci :
\add{si un groupe \emph{fini}}
qui \add{\uncertain{fidèle}} \struck{\ill{}} sur une carte, par
\struck{des} opérations isotopes à id, il opère trivialement
\uncertain{de même} : si un automorphisme d'ordre fini est
isotope à id, \struck{il \uncertain{opère} \ill{}} c'est id. Il suffit de
le voir pour un disque fermé et un pt sur
le \add{bord des} disques --- on voit d'abord que c'est l'identité
sur le bord, \struck{\ill{}} donc

\drawing{à droite du texte : un heptagone avec un point central marqué $s$ ; un disque contenant une petite boucle gribouillée attachée au bord ; plus bas, un disque avec une face découpée en polygones et un petit cercle ; tout en bas, un petit disque traversé d'un segment.}

\page{126}
\note{feuille de calculs et de croquis épars, sans texte suivi ; nous relevons les formules dans l'ordre approximatif de la page, de haut en bas et de gauche à droite. Plusieurs dessins (un personnage dans un triangle, un profil coiffé d'un chapeau) sont des griffonnages sans rapport apparent avec le calcul et ne sont que mentionnés.}

\drawing{une sphère vue en coupe, portant les points $0$ (en haut), $1$ (en bas), $\infty$ et $\tfrac12$ sur deux méridiens ; un triangle découpé en cases.}

\struck{$(0,1) \to \infty$} \qquad $(0,1) \to \infty$, \qquad $(\infty, 1) \to 0$ \qquad \struck{$\dfrac{az+b}{cz+d}$}

\[
  z \longmapsto \frac{z}{cz + (1-c)}
\]
$\dfrac{z}{cz+(1-c)} : \Bigl( \dfrac{cz}{cz+1-c} + {}$ \ill{}
\[
  \begin{pmatrix} 1 & 0 \\ c & 1-c \end{pmatrix}
  \begin{pmatrix} 1 & 0 \\ c & 1-c \end{pmatrix}
  = \begin{pmatrix} 1 & 0 \\ 2c - c^{2} & (1-c)^{2} \end{pmatrix}
\]
\note{le second facteur est en partie caché par le dessin d'une sphère ($0$, $1$, $\infty$ et $z$ marqués) ; nous le complétons d'après le produit écrit à droite.}
\[
  (1-c)^{2} = 1, \qquad c - 1 = \pm 1, \qquad c = 0 \text{ ou } 2 .
\]

\drawing{un pentagone finement triangulé autour d'un sommet central, les sommets $0$ et $1$ marqués au bord.}

\[
  \frac{z}{2z-1} = \frac{1}{2 - \frac{1}{z}}, \qquad
  z = \frac{z}{2z-1}, \qquad 2z^{2} - 2z = 0, \qquad z = 0 \text{ ou } 1 .
\]
\note{devant $2z^{2} - 2z = 0$, un premier $2z^{2}$ est biffé.}
\note{le dénominateur du membre de droite de la première égalité est lu $2 - \frac12$ ; le calcul exige $2 - \frac1z$.}

$g$ : $\underbrace{0,1,\infty}$ $\longrightarrow 0$ ; \quad $(0,$ \struck{$1$}$, \infty) \to (0, \infty)$ ; \quad
\struck{ou $0, 1$,} $1 \to (0,\infty)$ \quad \emph{ou} $1$.

$(0, 1, \infty) \longrightarrow (0, 1, \infty)$

\drawing{une sphère avec équateur et méridien en pointillé, les points $0$, $1$, $\infty$ marqués ; deux petits triangles portant le graphe dual (un « Y » intérieur), l'un avec le sommet $0$ ; un losange coupé par une diagonale et un « Y ».}

si \struck{\ill{} \ill{}} $s = 0, 1, \infty$ \uncertain{j'en} \ill{} \ill{} $1$
\ill{} \ill{} \ill{} \ill{} \ill{}

\drawing{un polygone dont le bord est doublé d'une courbe ondulée, avec un éventail de triangles issus d'un sommet intérieur ; à droite, plusieurs disques et ellipses avec un rayon marqué $z$, un cône.}

\[
  f(z) \in [0,1], \qquad f(z) - \tfrac12, \qquad \bigl(2f(z) - 1\bigr)^{2},
\]
\[
  g(z) = \bigl(2f(z) - 1\bigr)^{2},
\]
\struck{$g(z) = 0 \Rightarrow f(z) = $}
\[
  g(z) = 0 \iff f(z) = \tfrac12, \qquad g(z) = 1 \iff f(z) = 0, 1 .
\]
\note{les signes $\iff$ sont de nous : la page aligne les deux colonnes sans signe.}

\drawing{un segment $[0,1]$ avec son milieu marqué ; un pentagone aux sommets noircis et aux milieux d'arêtes marqués d'un trait.}

$\bigl(z - \tfrac12\bigr)^{2} \in [0,1]$, \quad \struck{$z = \tfrac{z}{2}$} \quad $z \in [0,1]$ ;
\struck{$z\bigl(z - \tfrac12\bigr) \in [\ldots$}
\[
  2z - 1 \in [-1, +1], \qquad (2z-1)^{2} \in [0, 1] .
\]

\page{127}
\note{cette page continue la phrase interrompue au bas de la page~123 (« car $\varphi^{-1}\{0,1,\infty\} \supset \{0,1,\infty\}$ et »), non la page~126.}
$= \varphi^{-1}\{0, \infty\} \supset \{0, 1, \infty\}$), avec indice de
ramification $2$ au-dessus de $1$ pour $\varphi f$.

\struck{\ill{}} Que se passe-t-il dans le cas exceptionnel ?

\page{128}
\struck{Soit} Soit $X$ une surface holomorphe, $C = (K, S)$ une
carte sur $X$. On dit que $C$ est \emph{compatible}
avec la structure holomorphe, ou encore
que $(X, K, S)$ est une \emph{carte holomorphe},
s'il existe une application holomorphe
$\varphi : X \longrightarrow \mathbb{P}^{1}_{\mathbb{C}} = S$, telle que
\begin{itemize}
  \item[a)] $\varphi$ est étale au-dessus de $S - \{0, 1, \infty\}$, et
  a le degré de ramification $2$ en tous les
  pts au-dessus de $1$.
  \item[b)] \struck{\ill{}} $(K, S) = \varphi^{-1}([0,1], 0)$ \note{un premier $K$ est biffé et récrit.}
\end{itemize}

\emph{Question} Unicité d'un tel $\varphi$ ? Soient $\varphi'$ ayant
les mêmes propriétés. Par le \uncertain{th.} de classification
il existe un unique isomorphisme de
\add{cartes, isotope à l'identité} \struck{holomorphe} rendant commutatif

\begin{tikzcd}
X \arrow[rr, "f"', "\simeq"] \arrow[dr, "\varphi"'] & & X \arrow[dl, "\varphi'"] \\
& S &
\end{tikzcd}

La question est donc si $f$ (qui est \uncertain{unique} \uncertain{holomorphe}\ldots) est l'identité, sachant
que c'est un automorphisme de la carte $C$
qui est isotope à id, i.e.\ fixant sommets,
arêtes, faces. \uncertain{Or} c'est faux déjà pour
la sphère \struck{$X$} $= S = \mathbb{P}^{1}_{\mathbb{C}}$, \uncertain{partagée} en deux
\drawing{en marge gauche, deux petites sphères esquissées, numérotées I) et II) (les cas exceptionnels de la page suivante).}

\page{129}
hémisphères par un équateur, avec $1$ ou
deux sommets dessus (dans \uncertain{chacun} des
cas, $1$ sommet et une arête, ou $2$ sommets
et $2$ arêtes). Le groupe des automorphismes
holomorphes qui fixe la carte et ses
él.\ de structure (donc induit un
automorphisme isotope à $1$) est le groupe des
$z \mapsto \dfrac{az+b}{cz+d}$ ($a, b, c, d \in \mathbb{R}$, $ad - bc = 1$) qui fixent
$\infty$ (donc $c = 0$ \add{i.e.\ le} groupe $\mathrm{Aff}(1, \mathbb{R})$ \struck{\ill{} \ill{}}),
\uncertain{soit} $0$ et $\infty$ (groupe $\mathbb{C}^{*}$ \struck{\ill{}} \uncertain{de} dim~$1$).
\note{« $\mathrm{Aff}(1, \mathbb{R})$ » : le second argument est surchargé ; nous lisons $\mathbb{R}$ d'après la condition $a, b, c, d \in \mathbb{R}$, mais la page porte peut-être $\mathbb{C}$ repassé.}
Autre exemple :

\drawing{un disque (sphère) traversé d'un arc joignant deux points, numéroté III).}

(1 arête, 2 sommets,
\uncertain{dual} de 1 arête 1 sommet) qui a
le $=$ groupe $\mathrm{Aff}(1, \mathbb{C})$. // Si on exige
que \struck{les \ill{} de} \struck{\add{chaque \uncertain{comp.} connexe de $X$ ait une}} \add{(\uncertain{chaque} face \add{\uncertain{deux} \uncertain{sommets}} \struck{soit} au}
moins trois, va-t-on \emph{gagner} ? J'ai
l'impression que oui, grâce à la
théorie de la représentation conforme.
\marginal{au moins \uncertain{pour} les cartes finies\ldots}
La condition envisagée signifie que
les composantes connexes des \uncertain{Uhe}\note{lecture incertaine ; peut-être une abréviation pour les faces} \uncertain{sont}
distinctes de la sphère munie
soit de la carte du type $(2, 1)$ (cas I)

\page{130}
soit les cartes d'\uncertain{oranges} d'ordre $n \geqslant 1$
--- mais si $n \geqslant 2$, il n'y a pas
d'automorphismes hol.\ $\neq \mathrm{id}$.

\drawing{une sphère découpée en fuseaux par des méridiens joignant deux pôles (« orange »).}

Donc il semble que les seuls cas
où $\varphi$ ne soit pas unique sont les cas
(isomorphes :) I), II), III) $=$ I$^{*}$) (\emph{NB} II) $=$ II$^{*}$)

Il faut cependant tirer au clair
entièrement l'\emph{argument} de représentation conforme, et voir s'il s'applique
aux cartes infinies.

Si on se donne en plus \struck{les} l'ens.\ \struck{\ill{}}
$\varphi^{-1}(\{\infty\})$ (centres des faces), alors les cas I) II) III)
sont éliminés. [Si on se donnait
$\varphi^{-1}(\{1\})$, cela éliminerait les cas
II et III, mais non le cas I\ldots]

\page{131}
\drawing{un cercle (la sphère $S$) portant le point $\infty$ et le segment de $0$ à $1$.}

\emph{Opération} \add{\uncertain{$\Omega$} orientée}, définie par une \struck{\ill{} complexe}
carte $(K', S')$ sur $S$, ayant éventuellement
un bout \uncertain{arrondi} en le pt $1 \in S$
\add{défin.\ \uncertain{à} une isotopie \uncertain{constante} sur $[0,1] \cup \{\infty\}$ près}.
Considérons les conditions
\begin{itemize}
  \item[a)] Pour toute carte holomorphe $C = (X, K, S)$, il existe
  \struck{\ill{} \ill{} $K'_X$ \ill{}} une représentante $(K'_C, S'_C)$ de $\Omega(C)$ qui est
  compatible avec la structure complexe.
  \item[b)] Pour toute \struck{\ill{} \ill{} dessus} \add{carte $(X, K, S)$}, les \uncertain{réalisations}
  \struck{complexes} \add{holomorphes} canoniques de celle-ci et de
  $(X, K'_C, S'_C)$ \note{le $X$ est surchargé.} \struck{sont} \uncertain{sont} \struck{\ill{}} des surfaces \marginal{\struck{b') comme b),}}
  \struck{\ill{}} holomorphes isomorphes.
    \item[c)] \struck{b)} Pour toute carte $C = (X, K, S)$ \add{et \uncertain{choix} de $K'_X$}, il existe
  une structure complexe sur $X$
  qui est compatible avec $(K, S)$, $(K'_X, S'_X)$ et
  \uncertain{avec} \uncertain{\ill{}} $\bigl((K,S), (K'_X, S'_X)\bigr)$.
\end{itemize}
a') b') c') comme a) b) c) mais \uncertain{avec} groupe d'autom.\ hol.
\begin{itemize}
  \item[d)] Il existe une application \emph{holomorphe}
  $\varphi : S \to S$, \struck{définissant} « spéciale » [i.e.\ \struck{ramifiée} \add{étale}
  au-dessus de $S - \{0, 1, \infty\}$,
  \struck{\ill{}} ayant ramification \struck{\ill{}} \add{$\leqslant$} $2$
  en \uncertain{les} pts au-dessus de $1$ \struck{\ill{}} \uncertain{\ill{}}~$1$,
  \struck{ou $1$}, \uncertain{envoyant} $\{0, \infty\}$ dans $\{0, \infty\}$ et
  $1$ en $\{0, \infty, 1\}$, et telle que \struck{\ill{}} $\varphi(1) = 1$
  \uncertain{impliqué} \uncertain{par} que $\varphi$ est étale en $1$)] telle que
  $(K', S') \underset{\text{isotope}}{\approx} \bigl[\varphi^{-1}([0,1]), \{0\}\bigr]$
  \add{\uncertain{rel.} \uncertain{à} $[0,1] \cup \{\infty\}$}.
\end{itemize}
\note{la page s'arrête sur cette condition ; la suite éventuelle n'est pas dans le lot.}

\drawing{en marge gauche, au niveau de d) : une parabole coupée par une droite en deux points, un segment à bout marqué ; en dessous, une barre en « T » renversé et une équerre.}

\end{document}
